\section{Conclusion}
\label{sec:conclusion}

We set out the finite-time no-catastrophe probability for a two-type population in which either type may trigger one global absorbing event. Because the intensity depends on the whole path, the natural object is an occupation-time transform of the two type abundances rather than a terminal-count generating function alone (\cref{sec:model}).

The calculation is exact because the backward system is triangular. Solving for $G$, shifting $S$ by its long-time root, linearising the Riccati equation, and changing variable to $z_2$ produces the Gauss hypergeometric equation and the boxed formula \cref{eq:S-final}. The type-2 catastrophe rate reshapes the autonomous coordinate seen by the whole system; the type-1 rate does not feed back into $G$. The limits $h$ and $r_{2,-}$ are forever-avoidance probabilities. Equal rates and zero-birth boundaries are recorded in \cref{app:limits}; the equation-level map to Antal and Krapivsky is \cref{app:AK-recovery}.

One natural reading is the early intracellular phase of \textit{Yersinia pestis}: early and adapted phenotypes, with catastrophe as irreversible loss of host-cell containment. The regimes of \cref{sec:biology} then compare orderings of $(\delta_1,\delta_2)$.

Beyond the work of this paper, one might consider the release distributions (\cref{sec:ultimate-release,fig:release-law}) more closely and their implications for a broad infection-success probability within the host. Analysis along these lines may be capable of contributing a quantitative justification for a rate hypothesis favouring MAT/GATE in which the second phase of \textit{Y.\ pestis} actively triggers burst of the containing cell.
